\subsection{The full length spectrum at the critical cusp scale}
\label{sec:critical-cluster}

The tight spectrum omits curves that become contractible when the cusps
are filled. At $n\asymp\sqrt g$, these curves contribute on the same
length scale as the nonseparating curves. Their contribution is not
Poisson: a two-cusp pair of pants contributes its whole length spectrum,
not only its boundary.

For $s>0$, let $P_s$ be the hyperbolic pair of pants with two cusps and
one geodesic boundary of length $s$. Write $K_s$ for its primitive
unoriented length counting measure, including the boundary once and
excluding parabolic classes. Equal lengths are counted with their
geometric multiplicities. Let
\begin{equation}
 \mu_c(\mathrm ds)=\frac{c^2}{8\pi^2}\sinh(s/2)\,\mathrm ds,
 \qquad \nu(\mathrm dt)=\frac{\cosh t-1}{t}\,\mathrm dt.
 \label{eq:cluster-intensities}
\end{equation}

\begin{theorem}[Critical primitive length spectrum]
\label{thm:critical-cluster}
Suppose $g\to\infty$ and $n/\sqrt g\to c\in(0,\infty)$, and let
$X$ have law $\mathbb P^{\mathrm{WP}}_{g,n}$. Then
\begin{equation}
 \sum_{\substack{\gamma\text{ primitive}\\\text{unoriented on }X}}
       \delta_{\ell_X(\gamma)}
 \ \Longrightarrow\ \Xi_c:=\mathcal N+\sum_{s\in\mathcal Q_c}K_s,
 \label{eq:critical-cluster}
\end{equation}
as locally finite counting measures on $[0,\infty)$ with the vague
topology. Here $\mathcal N$ and $\mathcal Q_c$ are independent Poisson
point processes of intensities $\nu$ and $\mu_c$. In particular, for
every nonnegative continuous compactly supported function $f$,
\begin{equation}
 \mathbb E\mathrm e^{-\langle f,\Xi_c\rangle}
 =\exp\left\{-\int(1-\mathrm e^{-f(t)})\,\nu(\mathrm dt)
 -\int(1-\mathrm e^{-\langle f,K_s\rangle})\,\mu_c(\mathrm ds)\right\}.
 \label{eq:cluster-laplace}
\end{equation}
\end{theorem}

Shen and Wu obtained the factorial-moment asymptotics for short
two-cusp pants and the limiting probability that such a curve occurs
at this cusp scale~\cite[Lemmas~28 and~33]{SW}. The corresponding Poisson mean is
$\mu_c([0,L])=c^2(\cosh(L/2)-1)/(4\pi^2)$ on their stated short
window. Theorem~\ref{thm:critical-cluster} concerns the full primitive
process on every fixed length window, including the correlated lengths
inside each pair of pants. The background intensity and the geometric
moment method come from Mirzakhani and Petri~\cite{MP}. The critical
regime needs the volume bounds of Mirzakhani--Zograf and
Shen--Wu~\cite{MZ,SW} together with the zero-density case of
Theorem~\ref{thm:volume} and Proposition~\ref{prop:bounded-boundaries}:
at $n\asymp\sqrt g$ the earlier estimates bound the genus ratio
$V_{g-1,n+2}/V_{g,n}$ only from above. The full uniform expansion is not
used.

The nonsimple part distinguishes this regime from closed surfaces.
He, Shen, Wu and Xue proved that the nonsimple systole of a
Weil--Petersson random closed surface is asymptotic to $\log g$,
with a more precise concentration about $\log g-\log\log g$~\cite[Theorem~1.1]{HSWX}.
Thus nonsimple curves leave every
fixed length window in that model. At the critical cusp scale here,
two-cusp pants instead retain nonsimple curves of bounded length.

\emph{Outline of the proof.} Short multicurves are organized by their cut
graphs. Using only the volume bounds of~\cite{MZ,SW},
Lemma~\ref{lem:critical-cuts} shows that in the limit a cut graph
contributes only if it has one component carrying the genus, possibly
with nonseparating loops, and two-cusp pairs of pants attached to it as
leaves. A bounded-multiplicity argument (Lemmas~\ref{lem:primitive-count}
and~\ref{lem:critical-complex}) reduces intersecting simple curves and
nonsimple curves to such cut graphs; hence, with probability tending to
one, every short nonsimple curve lies in a two-cusp pair of pants. The mixed
factorial moments of nonseparating lengths and pants waists converge to those
of independent Poisson processes (Lemma~\ref{lem:critical-skeleton}, where the
zero-density volume ratios enter), and attaching to each waist the full
spectrum of its pair of pants gives $\Xi_c$. The details are given in
Appendix~\ref{app:critical}.

The limit is not Poisson. If $m_L(s)=K_s([0,L])$, then
Lemma~\ref{lem:primitive-count} and the support restriction $s\le2L$
give finite moments and
\begin{equation}
 \operatorname{Var}\Xi_c([0,L])-\mathbb E\Xi_c([0,L])
 =\int m_L(s)(m_L(s)-1)\,\mu_c(\mathrm ds).
 \label{eq:cluster-overdispersion}
\end{equation}
Besides its waist, $P_s$ contains the primitive figure-eight around
its two cusps, of length $2\operatorname{arccosh}(2+\cosh(s/2))$
by \cite[Theorem~1.3]{Rivin}. For $L>2\operatorname{arccosh}3$,
both lengths lie below $L$ on an interval of positive $\mu_c$-measure.
The right side of \eqref{eq:cluster-overdispersion} is therefore
strictly positive. The theorem identifies these correlations in the
limit; no convergence of the full-spectrum second moments is asserted.
